% Ejemplo de tema: cada fichero de tema empieza con su \section.
% El número se muestra como «Tema 1» en el título y en el índice, y
% los entornos numeran por tema (Teorema 1.1, Definición 1.1, ...).

\section{Primer tema de la asignatura}\label{tema:uno}

Párrafo de introducción del tema. Los términos clave se anotan al
margen con \notamargen{término clave $\rightarrow$}, al estilo de los
libros de Diestel, para construir un mini-índice del tema.

\begin{definicion}[Objeto de estudio]\label{def:objeto}
    Definición del objeto de estudio del tema, con un nombre entre
    corchetes que aparece tras el número. Al estar el párrafo en
    cursiva y ser largo, se aprecia la sangría francesa: la primera
    línea queda al margen y las siguientes sangradas.
\end{definicion}

\begin{notacion}\label{not:basica}
    Notación básica del tema. El entorno de notación numera de forma
    independiente, igual que definiciones o ejemplos.
\end{notacion}

\begin{teorema}[Nombre del teorema]\label{thm:principal}
    Enunciado del teorema principal del tema, con la etiqueta opcional
    entre corchetes que aquí funciona de nombre propio y debajo, en
    \cref{eq:principal}, la ecuación de la que depende.
    \begin{equation}\label{eq:principal}
        a^2 + b^2 = c^2
    \end{equation}
\end{teorema}
\notamargen{dem. de \cref{thm:principal} usa \cref{def:objeto}}

\begin{demostracion}
    Demostración del teorema. La palabra «Demostración» está metida
    en la sangría francesa: abre la primera línea al margen y las
    líneas siguientes quedan sangradas; la demostración aparece
    siempre pegada a su teorema, sin salto de página posible entre
    ambos, y termina con el símbolo de QED.
\end{demostracion}

\begin{demostracion}[Prueba alternativa]
    Comentario: con un argumento opcional, la palabra «Demostración»
    se sustituye por el nombre que se indique, como aquí.
\end{demostracion}

\begin{lema}\label{lem:auxiliar}
    Enunciado del lema auxiliar, que comparte numeración con los
    teoremas y sus hermanos (proposiciones, corolarios, propiedades,
    axiomas y criterios).
\end{lema}

\begin{demostracion}
    Demostración del lema: por \cref{thm:principal} y la
    \cref{def:objeto}, con una referencia a la página \cpageref{def:objeto}.
\end{demostracion}

\begin{corolario}
    Enunciado del corolario del teorema principal.
\end{corolario}

\begin{proposicion}
    Enunciado de la proposición.
\end{proposicion}

\begin{propiedad}
    Enunciado de la propiedad.
\end{propiedad}

\begin{axioma}
    Enunciado del axioma.
\end{axioma}

\begin{criterio}
    Enunciado del criterio.
\end{criterio}

Un párrafo normal entre bloques, con una lista que también se
referencia con \cref{item:uno} e \cref{item:dos} y una nota al pie
referenciable\footnote{Contenido de la nota al pie, referenciable con
\cref{fn:nota}. \label{fn:nota}}.

\begin{enumerate}
    \item\label{item:uno} Primer ítem de la lista, referenciable como ítem.
    \item\label{item:dos} Segundo ítem, con \cref{item:uno} como ejemplo de referencia interna entre ítems.
\end{enumerate}

\begin{ejemplo}\label{ej:primero}
    Enunciado del ejemplo, que muestra cómo usar \cref{lem:auxiliar}.
\end{ejemplo}

\begin{observacion}
    Contenido de la observación.
\end{observacion}

\begin{nota}
    Contenido de la nota.
\end{nota}

\begin{conjetura}
    Enunciado de la conjetura que quedará por demostrar.
\end{conjetura}

\begin{problema}
    Enunciado del problema propuesto.
\end{problema}
